\documentclass[12pt,a4paper]{article}

\usepackage[a4paper,margin=18mm]{geometry}
\usepackage[T2A]{fontenc}
\usepackage[utf8]{inputenc}
\usepackage[russian]{babel}
\usepackage{amsmath,amssymb,mathtools}
\usepackage{array,booktabs}
\usepackage{xcolor}
\usepackage{hyperref}
\usepackage{tikz}

\hypersetup{
  colorlinks=true,
  linkcolor=blue!55!black,
  urlcolor=blue!55!black
}

\setlength{\parindent}{0pt}
\setlength{\parskip}{0.55em}
\setlength{\tabcolsep}{3pt}
\renewcommand{\arraystretch}{1.25}

\begin{document}

\begin{titlepage}
\centering
\vspace*{\fill}

{\LARGE\bfseries Ревью домашней работы\par}

\vspace{1.2cm}

{\large Автор: Лёвин Артём Александрович\par}

\vspace{0.6cm}

{\large 13 сентября 2026 года\par}

\vspace*{\fill}

\begin{tikzpicture}
\draw[line width=0.7pt] (0,0) .. controls (2,-0.25) and (4,0.25) .. (6,0);
\end{tikzpicture}

\end{titlepage}

\newpage

\section*{Сводная таблица}

{\small
\begin{center}
\begin{tabular}{p{0.11\linewidth} p{0.18\linewidth} p{0.28\linewidth} p{0.14\linewidth} p{0.16\linewidth}}
\toprule
\textbf{№ задания} & \textbf{Тема} & \textbf{Суть ошибки} & \textbf{Ваш ответ} & \textbf{Правильный ответ} \\
\midrule
\bottomrule
\end{tabular}
\end{center}
}

\vspace{1.5em}

\section*{Краткий комментарий}
По читаемой части работы ошибок, требующих отдельного разбора, не обнаружено. При оформлении фотографий важно, чтобы финальные пояснения и ответы полностью попадали в кадр: тогда решение можно проверить целиком и без догадок.

\vfill
\begin{center}
\small Лёвин Артём Александрович эксклюзивно для Володи Хачатуряна
\end{center}

\end{document}
